\documentclass[10pt,a4paper]{amsart}
\usepackage[margin=19mm]{geometry}
\usepackage{amsmath,amssymb,mathtools}
\usepackage[scr=rsfs,cal=euler]{mathalfa}
\usepackage{xfrac}
\usepackage[unicode,hidelinks,bookmarks=false]{hyperref}
\newcommand{\ie}{i.e.\ }
\newcommand{\cf}{cf.\ }
\newcommand{\resp}{respectively\ }
\newcommand{\Subsection}{\subsection}
\providecommand{\nobreakdash}{\nobreak\hbox{-}}
\setlength{\emergencystretch}{2em}
\multlinegap0pt
\usepackage{mathtools}
\usepackage{amssymb}
\usepackage[shortlabels]{enumitem}
\usepackage{xcolor}
\usepackage{cancel}
\usepackage{tikz}
\newtheorem{proposition}{Proposition}
\newcommand*\circled[1]{\tikz[baseline=(char.base)]{
            \node[shape=circle,draw,inner sep=2pt] (char) {#1};}}
\title{Functions with a maximal number of finite invariant or internally-1-quasi-invariant sets or supersets}
\author{Nizar El Idrissi, Samir Kabbaj}
\date{Source: arXiv:2210.11412v2 (CC BY 4.0)}
\begin{document}
\maketitle

\noindent\emph{Source and licence.} Functions with a maximal number of finite invariant or internally-1-quasi-invariant sets or supersets. Version arXiv:2210.11412v2 (CC BY 4.0), distributed by arXiv under the Creative Commons Attribution 4.0 International licence, http://creativecommons.org/licenses/by/4.0/, as recorded in the arXiv licence metadata for this exact version and shown on the abstract page https://arxiv.org/abs/2210.11412v2. Full licence notice: \texttt{licenses/arxiv-2210.11412-CC-BY-4.0.txt}.\par\smallskip

\noindent\textit{Editorial scope.} Counted unit: the article's proposition propBetaInt (source0.tex lines 531--542). The article's complete original proof of it is delivered with it (source0.tex lines 544--591, one \texttt{proof} environment of 48 lines). No mathematics is written, repaired or re-derived: the statement and the proof are the article's own lines, in the article's own notation, and no retained line is substituted or re-flowed.  No further source-local context is needed: every cross-reference of every retained line resolves inside the two delivered spans.  Nothing was omitted from any retained span, so the package carries no omission record.  No retained line contains a citation, so the wrapper carries no bibliography and no bibliography entry.  No retained line contains a figure environment, an image-inclusion command or drawing code, so no image is delivered and the package declares no asset dependency.  Editorial formatting only, in addition: an AMS \texttt{amsart} preamble that restates the article's own declarations and macros used by the retained lines, namely the environments \texttt{proposition} on separate counters, together with the standard packages those lines need.  The retained environments print in this document as Proposition 1 in order of appearance, whereas the article prints the same items with the numbering of its own full document.  The complete native source of the article as distributed in the arXiv e-print for this exact version is bundled unchanged as \texttt{sources/132366/source0.tex}.
\begin{proposition}
\label{propBetaInt}
Let $\beta : \mathbb{N} \to \mathbb{N}$ such that there exists a function $w : Int_{\omega,*}(\mathbb{N}) \to \mathbb{N}$ such that
\[ \forall [\![a,b]\!] \in Int_{\omega,*}(\mathbb{N}) : w([\![a,b]\!]) \in [\![a,b]\!] \text{ and } \beta([\![a,b]\!] \setminus \{w([\![a,b]\!])\}) \subseteq [\![a,b]\!]. \]
Let $\{b_n\}_{n \in [\![0,l]\!]}$ the sequence of non-fixed points of $\beta$ in strictly increasing order, where $l \in \mathbb{N} \cup \{-1, \infty\}$. We extend $\{b_n\}_{n \in [\![0,l]\!]}$ and $\beta$ by setting $b_{-1} = -1$ and $\beta(-1) = 0$, and if $l < \infty$, we set $b_{l+1} = \infty$. Then we have either \\
\textbf{Case \circled{1}}. $\forall n \in [\![0,l]\!] : b_n < \beta(b_n) \leq b_{n+1}$ and
\[ w = \begin{cases} Int_{\omega,*}(\mathbb{N}) &\to \mathbb{N} \\ [\![a,b]\!] &\mapsto b_n \text { if } \exists! n \in [\![0,l]\!] : b_n \in [\![a,b]\!] \text{ and } \beta(b_n) > b \\ [\![a,b]\!] &\mapsto \text{any element of } [\![a,b]\!] \text{ otherwise} \end{cases}. \]
\textbf{Case \circled{2}}. There exists $n^* \in [\![-1,l-1]\!]$ such that $\forall m \leq n^* : b_m < \beta(b_m) \leq b_{m+1}$, $\beta(b_{n^*+1}) < b_{n^*+1}$, $\forall m \in [\![n^*+2,l]\!] : b_{m-1} \leq \beta(b_m) < b_m$, and
\[ w = \begin{cases} Int_{\omega,*}(\mathbb{N}) &\to \mathbb{N} \\ [\![a,b]\!] &\mapsto b_n \text { if } \exists! n \in [\![0,n^*]\!] : b_n \in [\![a,b]\!] \text{ and } \beta(b_n) > b \\ [\![a,b]\!] &\mapsto b_n \text { if } \exists! n \in [\![n^*+1,l]\!] : b_n \in [\![a,b]\!] \text{ and } \beta(b_n) < a \\ [\![a,b]\!] &\mapsto \text{any element of } [\![a,b]\!] \text{ otherwise} \end{cases}. \]
\textbf{Case \circled{3}}. There exists $n^* \in [\![-1,l-1]\!]$ such that $\forall m \leq n^* : b_m < \beta(b_m) \leq b_{m+1}$, $\beta(b_{n^*+1}) > b_{n^*+2}$, $\forall m \in [\![n^*+2,l]\!] : b_{m-1} \leq \beta(b_m) < b_m$, and 
\[ w = \begin{cases} Int_{\omega,*}(\mathbb{N}) &\to \mathbb{N} \\ [\![a,b]\!] &\mapsto b_n \text { if } \exists! n \in [\![0,n^*]\!] : b_n \in [\![a,b]\!] \text{ and } \beta(b_n) > b \\ [\![a,b]\!] &\mapsto b_n \text { if } \exists! n \in [\![n^*+1,l]\!] : b_n \in [\![a,b]\!] \text{ and } \beta(b_n) < a \\ [\![a,b]\!] &\mapsto \text{any element of } [\![a,b]\!] \text{ otherwise} \end{cases}. \]
\end{proposition}

\begin{proof}
Let's show that 
\begin{equation}
\label{Eq1}
\boxed{\forall n \in [\![0,l]\!] : \left( \beta(b_n) < b_n \Rightarrow \forall m \in [\![n+1,l]\!] : b_{m-1} \leq \beta(b_m) < b_m \right) }
\end{equation}
and
\begin{equation}
\label{Eq2}
\boxed{\forall n \in [\![0,l]\!] : \left( \beta(b_n) > b_{n+1} \Rightarrow \forall m \in [\![n+1,l]\!] : b_{m-1} \leq \beta(b_m) < b_m \right) }.
\end{equation}
Suppose that $\beta(b_n) < b_n$ for a certain $n \in \mathbb{N}$ and $n+1 \leq l$. We have ${\beta([\![b_n,b_{n+1}]\!] \setminus \{w([\![b_n,b_{n+1}]\!])\}) \subseteq [\![b_n,b_{n+1}]\!]}$. If $w([\![b_n,b_{n+1}]\!]) = b_{n+1}$, then $b_n \leq \beta(b_n)$ which is a contradiction. Therefore $w([\![b_n,b_{n+1}]\!]) \neq b_{n+1}$ and so $b_n \leq \beta(b_{n+1}) < b_{n+1}$ since $b_{n+1}$ is not a fixed point. By immediate induction, we have ${\forall m \in [\![n+1,l]\!] : b_{m-1} \leq \beta(b_m) < b_m}$. \\
Therefore assertion \ref{Eq1} holds. \\
Suppose that $\beta(b_n) > b_{n+1}$ for a certain $n \in \mathbb{N}$ and $n + 1 \leq l$. We have ${\beta([\![b_n,b_{n+1}]\!] \setminus \{w([\![b_n,b_{n+1}]\!])\}) \subseteq [\![b_n,b_{n+1}]\!]}$. If $w([\![b_n,b_{n+1}]\!]) = b_{n+1}$, then ${\beta(b_n) \leq b_{n+1}}$ which is a contradiction. Therefore $w([\![b_n,b_{n+1}]\!]) \neq b_{n+1}$ and so ${b_n \leq \beta(b_{n+1}) < b_{n+1}}$ since $b_{n+1}$ is not a fixed point. By immediate induction, we have ${\forall m \in [\![n+1,l]\!] : b_{m-1} \leq \beta(b_m) < b_m}$. \\
Therefore assertion \ref{Eq2} holds. \\
With these assertions in hand, we can now start the proof. We treat the $l = \infty$ and $l < \infty$ cases simultaneously. \\
\textbf{Case \circled{1}}. Suppose $\forall n \in [\![0,l]\!] : b_n < \beta(b_n) \leq b_{n+1}$. \\
Let $a \leq b \in \mathbb{N}$. \\
\textbf{Subcase \circled{1.1}}. Suppose $\exists n \in [\![0,l]\!] : b_n \in [\![a,b]\!] \text{ and } \beta(b_n) > b$. Let's show that there is a unique such $n$. Suppose that $n_1 > n_2$ both satisfy the condition. Since, $n_2 > n_1$, we have $n_2 \geq n_1 + 1$, and since $(b_n)_{n \in \mathbb{N}}$ is strictly increasing, we have 
\[ b \geq b_{n_2} \geq b_{n_1 + 1} \geq \beta(b_{n_1}) > b, \]
which is a contradiction. Therefore there exists a unique $n(a,b) \in [\![0,l]\!]$ such that $b_{n(a,b)} \in [\![a,b]\!] \text{ and } \beta(b_{n(a,b)}) > b$. If we had $w([\![a,b]\!]) \neq b_{n(a,b)}$, then since $\beta([\![a,b]\!]\setminus \{w([\![a,b]\!])\}) \subseteq [\![a,b]\!]$, we would have $\beta(b_{n(a,b)}) \in [\![a,b]\!]$, which contradicts $\beta(b_{n(a,b)}) > b$. Therefore $w([\![a,b]\!]) = b_{n(a,b)}$. \\
\textbf{Subcase \circled{1.2}}. Suppose $\forall n \in [\![0,l]\!] : (b_n \in [\![a,b]\!] \Rightarrow \beta(b_n) \leq b)$. There is no additional constraint in this subcase, we can choose $w([\![a,b]\!])$ arbitrarily from $[\![a,b]\!]$. \\
\textbf{Conversely}. If $l \in \mathbb{N} \cup \{-1, \infty\}$, $\{b_n\}_{n \in [\![0,l]\!]}$ is a strictly increasing sequence, $\beta : \mathbb{N} \to \mathbb{N}$ is a function such that $\mathbb{N} \setminus \{b_n\}_{n \in [\![0,l]\!]}$ are the fixed points of $\beta$, $\forall n \in [\![0,l]\!] : b_n < \beta(b_n) \leq b_{n+1}$, and 
\[ w = \begin{cases} Int_{\omega,*}(\mathbb{N}) &\to \mathbb{N} \\ [\![a,b]\!] &\mapsto b_{n(a,b)} \text { if } \exists n \in [\![0,l]\!] : b_n \in [\![a,b]\!] \text{ and } \beta(b_n) > b \\ [\![a,b]\!] &\mapsto \text{any element of } [\![a,b]\!] \text{ otherwise} \end{cases}, \]
then $\forall a \leq b \in [\![a,b]\!] : \beta([\![a,b]\!] \setminus \{w([\![a,b]\!])\}) \subseteq [\![a,b]\!]$. Indeed, let $a \leq b \in \mathbb{N}$. Because of the uniqueness property of $b_{n(a,b)}$ (when it exists), any $x \in [\![a,b]\!] \setminus \{w([\![a,b]\!])\}$ is either a fixed point (in which case $\beta(x) = x \in [\![a,b]\!]$), or satisfies $x = b_n$ with $\beta(b_n) \leq b$, in which case, since $b_n < \beta(b_n)$, we have $a \leq x =  b_n < \beta(b_n) = \beta(x) \leq b$. \\

\noindent Suppose now $\neg \left( \forall n \in [\![0,l]\!] : b_n < \beta(b_n) \leq b_{n+1} \right)$. Necessarily, $l \geq 0$. Let $n^* = \max \{ n \in [\![-1,l]\!] : \forall m \leq n : b_m < \beta(b_m) \leq b_{m+1} \}$. Necessarily, $n^* \in [\![-1,l-1]\!]$. Notice that if $\beta(b_{n^* +1}) \leq b_{n^*+2}$, then by maximality of $n^*$, we have $\beta(b_{n^*+1}) < b_{n^*+1}$ since $b_{n^*+1}$ is not a fixed point of $\beta$. Therefore we have either $\beta(b_{n^*+1}) < b_{n^*+1}$ or $\beta(b_{n^*+1}) > b_{n^*+2}$. \\

\noindent \textbf{Case \circled{2}}.  Suppose $\forall m \leq n^* : b_m < \beta(b_m) \leq b_{m+1}$ and $\beta(b_{n^*+1}) < b_{n^*+1}$. By assertion \ref{Eq1}, we have $\forall m \in [\![n^*+2,l]\!] : b_{m-1} \leq \beta(b_m) < b_m$. \\
Let $a \leq b \in \mathbb{N}$. \\
\textbf{Subcase \circled{2.1}}. Suppose $\exists n \in [\![0,n^*]\!] : b_n \in [\![a,b]\!] \text{ and } \beta(b_n) > b$. Let's show that there is a unique such $n$. Suppose that $n_1 > n_2$ both satisfy the condition. Since, $n_2 > n_1$, we have $n_2 \geq n_1 + 1$, and since $(b_n)_{n \in \mathbb{N}}$ is strictly increasing, we have 
\[ b \geq b_{n_2} \geq b_{n_1 + 1} \geq \beta(b_{n_1}) > b, \]
which is a contradiction. Therefore there exists a unique $n(a,b) \in [\![0,n^*]\!]$ such that $b_{n(a,b)} \in [\![a,b]\!] \text{ and } \beta(b_{n(a,b)}) > b$. If we had $w([\![a,b]\!]) \neq b_{n(a,b)}$, then since $\beta([\![a,b]\!]\setminus \{w([\![a,b]\!])\}) \subseteq [\![a,b]\!]$, we would have $\beta(b_{n(a,b)}) \in [\![a,b]\!]$, which contradicts $\beta(b_{n(a,b)}) > b$. Therefore $w([\![a,b]\!]) = b_{n(a,b)}$. \\
\textbf{Subcase \circled{2.2}}. Suppose $\exists n \in [\![n^*+1,l]\!] : b_n \in [\![a,b]\!] \text{ and } \beta(b_n) < a $. Let's show that there is a unique such $n$. Suppose that $n_1 > n_2$ both satisfy the condition. Since, $n_2 > n_1$, we have $n_2 -  1 \geq n_1$ and $n_2 \geq n^*+2$, and since $(b_n)_{n \in \mathbb{N}}$ is strictly increasing, we have 
\[ a > \beta(b_{n_2}) \geq b_{n_2 - 1} \geq b_{n_1} \geq a, \]
which is a contradiction. Therefore there exists a unique $n'(a,b) \in [\![n^*+1,l]\!]$ such that $b_{n'(a,b)} \in [\![a,b]\!] \text{ and } \beta(b_{n'(a,b)}) > b$. If we had $w([\![a,b]\!]) \neq b_{n'(a,b)}$, then since $\beta([\![a,b]\!]\setminus \{w([\![a,b]\!])\}) \subseteq [\![a,b]\!]$, we would have $\beta(b_{n'(a,b)}) \in [\![a,b]\!]$, which contradicts $\beta(b_{n'(a,b)}) > b$. Therefore $w([\![a,b]\!]) = b_{n'(a,b)}$. \\
Let's show that subcases 2.1 and 2.1 are disjoint. Suppose the contrary. Then 
\begin{align*}
a > \beta(b_{n'(a,b)}) &\geq b_{n'(a,b)-1}  > \beta(b_{n'(a,b)-1}) \geq b_{n'(a,b)-2} > \cdots  \\
&> \beta(b_{n^*+2}) \geq b_{n^* + 1} \geq \beta(b_{n^*}) > b_{n^*} \geq \beta(b_{n^* - 1}) > \cdots > \beta(b_{n(a,b)}) > b, 
\end{align*}
which is a contradiction. \\
\textbf{Subcase \circled{2.3}}. Suppose $\forall n \in [\![0,n^*]\!] : (b_n \in [\![a,b]\!] \Rightarrow \beta(b_n) \leq b)$ and ${\forall n \in [\![n^*+1,l]\!] : (b_n \in [\![a,b]\!] \Rightarrow \beta(b_n) \geq a)}$. There is no additional constraint in this subcase, we can choose $w([\![a,b]\!])$ arbitrarily from $[\![a,b]\!]$. \\
\textbf{Conversely}. If $l \geq 0$, $\{b_n\}_{n \in [\![0,l]\!]}$ is a strictly increasing sequence, $\beta : \mathbb{N} \to \mathbb{N}$ is a function such that $\mathbb{N} \setminus \{b_n\}_{n \in [\![0,l]\!]}$ are the fixed points of $\beta$, $n^* \in [\![-1,l-1]\!]$, $\forall m \leq n^* : b_m < \beta(b_m) \leq b_{m+1}$, $\beta(b_{n^*+1}) < b_{n^*+1}$, $\forall m \in [\![n^*+2,l]\!] : b_{m-1} \leq \beta(b_m) < b_m$, and 
\[ w = \begin{cases} Int_{\omega,*}(\mathbb{N}) &\to \mathbb{N} \\ [\![a,b]\!] &\mapsto b_{n(a,b)} \text { if } \exists n \in [\![0,n^*]\!] : b_n \in [\![a,b]\!] \text{ and } \beta(b_n) > b \\ [\![a,b]\!] &\mapsto b_{n'(a,b)} \text { if } \exists n \in [\![n^*+1,l]\!] : b_n \in [\![a,b]\!] \text{ and } \beta(b_n) < a \\ [\![a,b]\!] &\mapsto \text{any element of } [\![a,b]\!] \text{ otherwise} \end{cases}, \]
then $\forall a \leq b \in [\![a,b]\!] : \beta([\![a,b]\!] \setminus \{w([\![a,b]\!])\}) \subseteq [\![a,b]\!]$. Indeed, let $a \leq b \in \mathbb{N}$. Because of the uniqueness properties of $b_{n(a,b)}$ and $b_{n'(a,b)}$ (when they exist) and the disjointness of subcases 2.1 and 2.2, any $x \in [\![a,b]\!] \setminus \{w([\![a,b]\!])\}$ is either a fixed point (in which case $\beta(x) = x \in [\![a,b]\!]$), or satisfies $x = b_n$ with $n \in [\![0,n^*]\!]$ and $\beta(b_n) \leq b$, in which case, since $b_n < \beta(b_n)$, we have $a \leq x =  b_n < \beta(b_n) = \beta(x) \leq b$, or satisfies $x = b_n$ with $n \geq n^* + 1$ and $\beta(b_n) \geq a$, in which case, since $\beta(b_n) < b_n$, we have $a \leq \beta(b_n) = \beta(x) < b_n = x \leq b$.  \\
\textbf{Case \circled{3}}.  Suppose $\forall m \leq n^* : b_m < \beta(b_m) \leq b_{m+1}$ and $\beta(b_{n^*+1}) > b_{n^*+2}$. By assertion \ref{Eq2}, we have $\forall m \geq n^* + 2 : b_{m-1} \leq \beta(b_m) < b_m$. Then the treatment from now on is analogous to case 2 (since we haven't used the conditions $\beta(b_{n^*+1}) < b_{n^*+1}$ or $\beta(b_{n^*+1}) > b_{n^*+2}$) with the same necessary and sufficient form of the function $w$.\\
\end{proof}

\end{document}
